% Generated draft: requires mathematical and source review.
\documentclass[11pt]{article}
\usepackage[margin=25mm]{geometry}
\usepackage{fontspec}
\setmainfont{Latin Modern Roman}
\usepackage{amsmath,amssymb,mathtools,unicode-math}
\setmathfont{Latin Modern Math}
\usepackage{enumitem,needspace,xurl,hyperref}
\setlength{\parindent}{0pt}
\setlength{\parskip}{5pt}
\setlength{\emergencystretch}{3em}
\newcommand{\N}{\mathbb{N}}
\newcommand{\Z}{\mathbb{Z}}
\newcommand{\Q}{\mathbb{Q}}
\newcommand{\R}{\mathbb{R}}
\newcommand{\C}{\mathbb{C}}
\newcommand{\F}{\mathbb{F}}
\newcommand{\A}{\mathbb{A}}
\newcommand{\PP}{\mathbb{P}}
\newcommand{\E}{\mathbb{E}}
\newcommand{\eps}{\varepsilon}
\newcommand{\dd}{\,\mathrm{d}}
\begin{document}
\section*{UnsolvedMath ID: 700803}
\textbf{Source problem code:} catalogue entry
\subsection{Definitions and mathematical statement}
A group \(G\) is \emph{locally finite} if every subgroup generated by finitely many elements of \(G\) is finite. A group \(S\) is \emph{simple} if \(S\neq\{1\}\) and its only normal subgroups are \(\{1\}\) and \(S\).

\textbf{First question.} Is it true that for every infinite locally finite simple group \(G\), there exists a subgroup \(S\leq G\) such that \(S\) is finite, nonabelian, and simple?

\textbf{Second question.} Does every infinite locally finite simple group \(G\) contain an ``infinite tower'' of nonabelian finite simple subgroups? Interpreting an infinite tower as a strictly ascending sequence, the question becomes whether, for every such \(G\), there exist subgroups \((S_n)_{n\geq 1}\) satisfying
\[
S_1<S_2<S_3<\cdots\leq G,
\qquad
\text{each }S_n\text{ is finite, nonabelian, and simple},
\]
where \(S_n<S_{n+1}\) means proper subgroup inclusion. The supplied source does not explicitly define ``infinite tower''; the strictly ascending interpretation is therefore an unresolved qualification. No requirement that the union of the tower equal \(G\) is stated.
\subsection{Short English statement}
B. Hartley asks whether every infinite locally finite simple group contains a nonabelian finite simple subgroup, and whether it contains an infinite tower of such subgroups. The meaning of ``tower'' requires confirmation; a strictly ascending sequence is the interpretation displayed here.
\subsection{Sources}
\begin{itemize}
\item Local UnsolvedMath registry, permanent ID \(700803\), exact source problem code \texttt{catalogue entry}, title \emph{Kourovka Notebook Problem 11.111}. The supplied registry labels the problem \texttt{open}; this label is source metadata, not an independently verified mathematical or literature assessment.
\item \emph{The Kourovka Notebook (Unsolved Problems in Group Theory)}, original problem 11.111, 11th issue (1990), proposed by B. Hartley. The supplied record also identifies the 21st issue edition (2026), \texttt{arXiv:1401.0300v48}. The collection description attributes editorship to E. I. Khukhro and V. D. Mazurov. These bibliographic details are supplied attribution and have not been independently verified.
\end{itemize}
\end{document}
